\section{Related Work}

The literature review used for this project identifies four main directions relevant to MOF adsorption benchmarking:
classical GCMC simulations for CO2 capture, active learning for mixture adsorption landscapes, machine-learned interatomic potentials for adsorption Monte Carlo, and flexible-framework workflows coupling molecular dynamics with sorption simulations.
These directions motivate the modular structure of \texttt{MOF\_benchwork} \cite{tao2022gcmc,mukherjee2023active,edwards2026mlipmc,wang2026flexible}.

\subsection{GCMC Simulations for Gas Adsorption in MOFs}

Grand Canonical Monte Carlo is a standard method for computing adsorption isotherms because it samples particle insertion and deletion at fixed temperature and chemical potential \cite{tao2022gcmc,mukherjee2023active,edwards2026mlipmc}.
For adsorption benchmarks, this is especially useful because the simulation directly produces pressure-dependent uptake values that can be compared with reference isotherms.

Tao et al. provide the closest methodological starting point for the current implementation \cite{tao2022gcmc}.
Their study investigates pure CO2 adsorption and CO2/N2 mixtures in MOF-5, Mg-MOF-74, and ZIF-8.
The work is relevant because it combines a classical GCMC setup, explicit pressure ranges, fugacity correction via the Peng-Robinson equation of state, long-range electrostatics, and validation against experimental data.
It also reports that GCMC yields absolute adsorption, while experiments often report excess adsorption.
This directly motivates the current implementation of both absolute and excess loading in the evaluation pipeline.

\subsection{MOF-5 / IRMOF-1 as First Benchmark System}

MOF-5, also known as IRMOF-1, is a suitable first benchmark material because it is a well-known rigid framework and appears in literature studies of CO2 adsorption \cite{tao2022gcmc,crafted2023}.
Compared with more complex systems such as open-metal-site MOFs or flexible water-responsive MOFs, MOF-5 is a practical baseline for validating the technical pipeline.

The current project therefore starts with CO2 in MOF-5/IRMOF-1 as Module A.
The goal is not yet to solve every adsorption problem, but to establish a reproducible path from configuration to LAMMPS input generation, execution, parsing, and reference comparison.

\subsection{Force Fields and Charge Models}

Classical force fields are computationally efficient and make long Monte Carlo simulations feasible.
However, the reviewed literature also shows that their reliability depends strongly on the material and adsorbate.
UFF or DREIDING-style force fields may be sufficient for initial rigid-framework benchmarks, but they can become unreliable for open metal sites, strongly polar adsorbates, cooperative interactions, or flexible frameworks \cite{tao2022gcmc,edwards2026mlipmc,wang2026flexible}.

This motivates the current design decision to keep force-field resolution explicit and reproducible.
For Module A, the framework is treated as rigid, framework charges are resolved from the selected CIF charge scheme, and the CO2 adsorbate model is read from CRAFTED/RASPA-style definition files.
This keeps the first benchmark simple enough to debug while preserving the information needed for later force-field comparisons \cite{crafted2023,edwards2026mlipmc}.

\subsection{Reference Data and Adsorption Basis}

Reference data are not automatically comparable just because they describe the same adsorbent and adsorbate.
Relevant differences include temperature, pressure grid, material naming, charge scheme, force field, and adsorption basis.
The adsorption basis is particularly important: simulation output from GCMC molecule counts is naturally absolute adsorption, while experimental sources may report excess adsorption \cite{tao2022gcmc,nistadsorption}.

The current project therefore records the reference basis where possible and evaluates references by basis:
absolute or simulation-reference data are compared against absolute loading, excess references are compared against excess loading, and unknown-basis references are explicitly marked.
CRAFTED is used as a structured simulation-reference source, while NIST ISODB is used as an additional literature and experimental reference source \cite{crafted2023,craftedzenodo,nistadsorption}.

\subsection{Mixture Adsorption and Active Learning}

Mixture adsorption introduces a larger thermodynamic feature space because pressure, composition, and temperature can vary simultaneously.
Mukherjee et al. show that active learning can reduce the number of required GCMC simulations for multicomponent adsorption landscapes \cite{mukherjee2023active}.
Their work uses GCMC data as reference data and applies Gaussian process regression to select new simulation points based on model uncertainty.

This motivates Module B as a future extension.
After the single-component baseline is validated, the benchmark can be extended to CO2/N2 mixtures and later to reduced sampling strategies.
In that stage, relevant metrics will include component uptake, selectivity, mean relative error, coefficient of determination, and number of GCMC points saved \cite{tao2022gcmc,mukherjee2023active}.

\subsection{Machine-Learned Interatomic Potentials}

Edwards et al. address the question of potential quality in adsorption Monte Carlo \cite{edwards2026mlipmc}.
Their MLIP-MC framework enables Widom insertion and GCMC calculations using machine-learned interatomic potentials.
The central conclusion from the literature review is that universal MLIPs are not automatically reliable for quantitative adsorption predictions.
Some models overbind and fill pores too strongly, while others underpredict adsorption.

This motivates Module C as a potential benchmark rather than an assumption that MLIPs are always better.
The project should eventually compare classical force fields, universal MLIPs, adsorption-relevant MLIPs, and reference models using quantities such as Widom insertion energy, Henry coefficients, isosteric heat of adsorption, isotherm error, and runtime \cite{edwards2026mlipmc}.

\subsection{Framework Flexibility}

Rigid-framework GCMC is not sufficient for every MOF.
Wang et al. show that water adsorption in NbOFFIVE-1-Ni requires a flexible-framework treatment because guest molecules affect the framework structure \cite{wang2026flexible}.
Their workflow couples MLIP-based molecular dynamics with GCMC sorption steps, allowing the framework to relax in response to adsorbate loading.

This motivates Module D as an advanced extension.
It should come after the rigid-framework pipeline is validated because it requires additional infrastructure for structural states, molecular dynamics, MLIP models, trajectory analysis, and convergence across sorption-relaxation cycles \cite{wang2026flexible}.

\subsection{Research Gap and Project Motivation}

The reviewed literature shows that many simulation methods and datasets already exist, but reproducible benchmarking remains difficult.
A practical benchmark must specify the material, structure preparation, charge model, force field, adsorbate model, fugacity or chemical potential conversion, simulation length, pressure range, reference data, adsorption basis, and uncertainty treatment \cite{tao2022gcmc,mukherjee2023active,edwards2026mlipmc,wang2026flexible,crafted2023,nistadsorption}.

\texttt{MOF\_benchwork} addresses this gap by turning a benchmark JSON file into a resolved run plan, generated LAMMPS inputs, parsed results, and curated reference comparisons.
The current Module A implementation is therefore the first reproducible baseline from which mixture, potential, and flexibility benchmarks can be developed.
